\documentclass[11pt,twocolumn]{article}

% ── Margins ────────────────────────────────────────────────────────────────
\usepackage[top=1in,bottom=1in,left=1in,right=1in]{geometry}

% ── Packages ───────────────────────────────────────────────────────────────
\usepackage{times}
\usepackage{graphicx}
\usepackage{booktabs}
\usepackage{amsmath}
\usepackage{hyperref}
\usepackage{enumitem}
\usepackage{microtype}
\usepackage{caption}
\usepackage{subcaption}
\usepackage[T1]{fontenc}

\setlength{\columnsep}{0.25in}
\captionsetup{font=small, labelfont=bf}
\setlist[itemize]{noitemsep, topsep=2pt, leftmargin=*}

% ── Title ──────────────────────────────────────────────────────────────────
\title{\textbf{NeuroScribe: EEG Seizure Detection with\\
Evidence-Grounded Clinical Reporting}\\[4pt]
\large Checkpoint 2 Report}
\author{cm7513}
\date{}

\begin{document}
\maketitle
\thispagestyle{empty}

% ══════════════════════════════════════════════════════════════════════════
\section{Dataset \& Preprocessing}
% ══════════════════════════════════════════════════════════════════════════

We use the \textbf{CHB-MIT Scalp EEG Database}~\cite{shoeb2009}, a publicly
available dataset of 24 pediatric patients with intractable epilepsy recorded
at Boston Children's Hospital at 256\,Hz with expert-annotated seizure times.
Signal preprocessing follows standard EEG analysis protocols~\cite{niedermeyer2004}.

\subsection{Data Statistics}

\begin{table}[h]
\centering
\caption{Dataset Statistics}
\small
\begin{tabular}{@{}ll@{}}
\toprule
\textbf{Property} & \textbf{Value} \\
\midrule
Patients           & 24 \\
EDF recordings     & 686 \\
EEG channels       & 23 bipolar \\
Sampling rate      & 256 Hz \\
Window size        & 4\,s (1024 samples) \\
Feature dimension  & $23 \times 1024$ \\
Total seizures     & 193 \\
Median seizure dur.& 46\,s \\
Signal range       & $[-721,\,+595]$\,$\mu$V \\
\bottomrule
\end{tabular}
\label{tab:stats}
\end{table}

\vspace{-6pt}
\subsection{Preprocessing Pipeline}

The following steps are applied sequentially in \texttt{src/data/dataset.py}:
\begin{itemize}
  \item \textbf{Channel standardization:} pin all EDFs to 23 common bipolar
        channels; skip files with incompatible layouts (chb12, chb13, chb16).
  \item \textbf{Bandpass filter} (0.5--40\,Hz): removes DC drift and
        high-frequency muscle artefacts.
  \item \textbf{Notch filter} (60\,Hz): removes US powerline interference.
  \item \textbf{Windowing:} 4-second non-overlapping windows.
  \item \textbf{Z-score normalization:} per channel per window
        (std: 40.64\,$\mu$V $\to$ 1.0).
\end{itemize}

\subsection{Data Split Strategy}

We use a strict \textbf{patient-independent split} — no patient appears in
more than one partition — to simulate real clinical deployment on unseen
subjects.

\begin{table}[h]
\centering
\caption{Patient-Independent Data Splits}
\small
\begin{tabular}{@{}lllll@{}}
\toprule
\textbf{Split} & \textbf{Patients} & \textbf{Windows} &
\textbf{Seizure} & \textbf{\%} \\
\midrule
Train & chb01--18 (18) & 704,679 & 3,842 & 0.55\% \\
Val   & chb19--21 (3)  & 80,421  & 183   & 0.23\% \\
Test  & chb22--24 (3)  & 70,968  & 257   & 0.36\% \\
\bottomrule
\end{tabular}
\label{tab:splits}
\end{table}

The severe class imbalance (1:200) is addressed via \textbf{stratified
subsampling}: all seizure windows are retained; non-seizure windows are
subsampled to a 1:5 ratio (16.67\% seizure rate), giving 23,052 / 1,098 /
1,542 windows for train/val/test respectively.

\subsection{Data Visualizations}

\begin{figure}[h]
  \centering
  \includegraphics[width=\columnwidth]{report_figs/seizure_waveform.png}
  \caption{Full EEG recording (chb01\_03.edf) — first 6 channels shown.
           Red shading marks the annotated seizure window (2996--3036\,s).}
  \label{fig:waveform}
\end{figure}

\begin{figure}[h]
  \centering
  \includegraphics[width=\columnwidth]{report_figs/class_distribution.png}
  \caption{Class distribution per split. Seizure windows are rare
           ($<$1\%) across all splits, motivating Focal Loss and
           stratified subsampling.}
  \label{fig:class_dist}
\end{figure}

% ══════════════════════════════════════════════════════════════════════════
\section{Baseline Results}
% ══════════════════════════════════════════════════════════════════════════

\subsection{CNN+GRU Seizure Detector}

The baseline model is a \textbf{1D CNN + Bidirectional GRU + Attention}
classifier~\cite{cho2014gru} implemented in \texttt{src/models/cnn\_gru.py}.
The CNN stage consists of 3 layers with 32$\to$64$\to$128 filters
(kernels 7/5/3), each followed by BatchNorm, ReLU, and MaxPool.
The temporal features are then passed to a 2-layer bidirectional GRU
with hidden size 128, producing a 256-dim output. A soft attention
mechanism over all time steps focuses the model on seizure-relevant
moments within each window. Dropout (0.5) is applied throughout,
yielding $\sim$650K trainable parameters.

\textbf{Training:} Focal Loss ($\alpha{=}0.85$, $\gamma{=}2.0$), Adam
(lr$=3\times10^{-4}$, wd$=10^{-4}$), gradient clipping (max norm$=1.0$),
batch size 512, early stopping (patience$=10$).

\begin{figure}[h]
  \centering
  \includegraphics[width=\columnwidth]{report_figs/training_curves.png}
  \caption{Training and validation loss/F1 curves over 22 epochs.
           Best model saved at epoch 12 (Val F1 = 0.577).}
  \label{fig:curves}
\end{figure}

\subsection{Test Set Results}

The best checkpoint (epoch 12) was evaluated on held-out patients
chb22--24 at threshold $\tau{=}0.45$ (tuned on validation set):

\begin{table}[h]
\centering
\caption{Final Test Set Results — CNN+GRU Baseline}
\small
\begin{tabular}{@{}ll@{}}
\toprule
\textbf{Metric} & \textbf{Value} \\
\midrule
Sensitivity (Recall) & 0.7237 \\
Precision            & 0.4869 \\
Specificity          & 0.8475 \\
F1 Score             & \textbf{0.5822} \\
ROC AUC              & \textbf{0.8591} \\
PR AUC               & 0.6467 \\
FP / hour            & \textbf{0.22} \\
\bottomrule
\end{tabular}
\label{tab:results}
\end{table}

\begin{table}[h]
\centering
\caption{Per-Patient Test Results}
\small
\begin{tabular}{@{}lllll@{}}
\toprule
\textbf{Patient} & \textbf{Sens.} & \textbf{Prec.} &
\textbf{F1} & \textbf{FP/hr} \\
\midrule
chb22 & 0.700 & 0.417 & 0.522 & 0.05 \\
chb23 & 0.926 & 0.521 & 0.667 & 0.10 \\
chb24 & 0.515 & 0.481 & 0.498 & 0.06 \\
\bottomrule
\end{tabular}
\label{tab:perpatient}
\end{table}

\begin{figure}[h]
  \centering
  \includegraphics[width=\columnwidth]{report_figs/confusion_roc.png}
  \caption{Confusion matrix (left) and ROC curve (right) on test set.
           ROC AUC = 0.896 confirms strong discriminative ability.}
  \label{fig:roc}
\end{figure}


% ══════════════════════════════════════════════════════════════════════════
\section{Codebase}
% ══════════════════════════════════════════════════════════════════════════

All source code is organized under \texttt{src/} with clear separation:

\begin{table}[h]
\centering
\caption{Source Code Structure}
\scriptsize
\begin{tabular}{@{}p{3.2cm}p{3.8cm}@{}}
\toprule
\textbf{File} & \textbf{Role} \\
\midrule
\texttt{data/loader.py}      & CHB-MIT parser, EDF loader \\
\texttt{data/dataset.py}     & Preprocessing, windowing, DataLoaders \\
\texttt{models/cnn\_gru.py}  & CNN+BiGRU+Attention classifier \\
\texttt{models/losses.py}    & Focal Loss \\
\texttt{training/trainer.py} & run\_epoch() + gradient clipping \\
\texttt{features/extractor.py} & Welch PSD, RMS, spatial features \\
\texttt{llm/report\_generator.py} & GPT-4o-mini report generator \\
\texttt{llm/claim\_verifier.py}   & Hallucination detection \\
\bottomrule
\end{tabular}
\label{tab:code}
\end{table}

Three notebooks demonstrate the full pipeline end-to-end:
\texttt{01\_eda\_preprocessing.ipynb},
\texttt{02\_baseline\_cnn\_gru.ipynb},
\texttt{03\_hallucination\_demo.ipynb}.

% ══════════════════════════════════════════════════════════════════════════
\section{Responses to Reviewer Questions}
% ══════════════════════════════════════════════════════════════════════════

\textbf{Q: LLM report generation needs further elaboration — suggest using a
trusted LLM to verify another LLM's report.}
In the final NeuroScribe system, we adopt a two-stage LLM pipeline inspired
by LLM-as-a-Judge frameworks~\cite{zheng2023judging}: a lightweight
\emph{Report Generator} (GPT-4o-mini) produces the initial clinical report
from extracted signal features, while a more capable \emph{Evidence Verifier}
(GPT-4o) independently audits every quantitative claim against the ground-truth
feature dictionary, flagging hallucinated values and correcting them before
the report is delivered. This follows the principle that a stronger model can
reliably evaluate and improve the output of a weaker one~\cite{zheng2023judging}.

\textbf{Q: The ground truth for hallucination evaluation is unclear.}
Ground truth has two fully objective components: (1) seizure onset/offset
times from expert neurologist annotations provided with the CHB-MIT dataset —
the accepted benchmark used in hundreds of published papers since 2010;
(2) signal features computed deterministically via standard signal processing:
Welch's PSD for dominant frequency, RMS for amplitude, per-channel ranking
for spatial features. These computations are exact and reproducible.

\textbf{Q: The 5-stage pipeline may lack necessary method designs.}
Each stage has explicit design choices: Focal Loss ($\alpha{=}0.85$,
$\gamma{=}2.0$) and stratified sampling address the 1:200 class imbalance;
soft attention focuses on seizure-relevant time steps; gradient clipping
(max norm$=1.0$) stabilizes training; threshold tuning on the validation set
optimizes the precision-sensitivity trade-off. The prompting strategy uses
structured feature injection with explicit constraints, reducing hallucination
from 83\% to $\sim$0\%.

\textbf{Q: Can GPT-4o-mini utilize extracted features without fine-tuning?}
Yes — GPT-4o-mini supports in-context learning natively. Fine-tuning is
unnecessary because we are not asking the model to learn new knowledge; we
provide the correct values directly in the prompt. The 83\%~$\to$~$\sim$0\%
hallucination reduction empirically confirms the model successfully uses the
provided features.

% ══════════════════════════════════════════════════════════════════════════
\begin{thebibliography}{9}
\bibitem{shoeb2009}
A.~H. Shoeb, ``Application of machine learning to epileptic seizure onset
detection and treatment,'' Ph.D. thesis, MIT, 2009. CHB-MIT Scalp EEG
Database, PhysioNet. \url{https://doi.org/10.13026/C2K01R}

\bibitem{niedermeyer2004}
E.~Niedermeyer and F.~H. Lopes da Silva, \textit{Electroencephalography:
Basic Principles, Clinical Applications, and Related Fields}, 5th ed.
Lippincott Williams \& Wilkins, 2004.

\bibitem{cho2014gru}
K.~Cho et al., ``Learning Phrase Representations using RNN Encoder-Decoder
for Statistical Machine Translation,''
\textit{Proceedings of EMNLP}, 2014.

\bibitem{zheng2023judging}
L.~Zheng et al., ``Judging LLM-as-a-Judge with MT-Bench and Chatbot Arena,''
\textit{Advances in Neural Information Processing Systems (NeurIPS)}, 2023.
\end{thebibliography}

\end{document}
